\documentclass[11pt]{article}
\usepackage[margin=1in]{geometry}
\usepackage{booktabs}
\usepackage{array}
\usepackage{enumitem}
\usepackage[htt]{hyphenat}
\sloppy
\usepackage[colorlinks=true,urlcolor=blue,linkcolor=black]{hyperref}

\title{Wearables and Habits: Setup, Data Download and Package Audit\\\large Issue \#3}
\date{3 October 2026}
\author{}

\begin{document}
\maketitle

\section{Summary}
The existing research container builds and runs. Open e-commerce 1.0 (CC0) is downloaded and checksum-verified. The
dataset matches the plan's description: Amazon purchases linked to a survey by respondent ID. The staggered-DiD stack
(R \texttt{did}, \texttt{fixest}, \texttt{HonestDiD}, \texttt{didimputation}, \texttt{bacondecomp}; Python \texttt{pyfixest})
is added to the same container. No new container was needed. This is a setup report; the full audit and go/no-go are the next issue.

\section{Data}
Source: Harvard Dataverse, DOI \href{https://doi.org/10.7910/DVN/YGLYDY}{10.7910/DVN/YGLYDY}, license CC0 1.0.
Downloaded by \texttt{code/scripts/download\_open\_ecommerce.py} (retry/backoff, HTTP range resume, md5 check) to
\texttt{code/build/input/public/open\_ecommerce/} (gitignored). All md5 sums match Dataverse.

\begin{center}\small
\begin{tabular}{p{5.4cm}rp{6.2cm}}
\toprule
File & Size & Content \\\midrule
\texttt{amazon-purchases.csv} & 313 MB & Order date, unit price, quantity, ship-to state, title, ASIN, category, respondent ID \\
\texttt{survey.csv} & 1.3 MB & One row per respondent: age group, race, education, income, gender, state, household size, order frequency, smoking/alcohol/diabetes/wheelchair, 2021 life changes, data-sharing attitudes \\
\texttt{fields.csv} & 2.5 kB & Question text for survey columns \\
\texttt{survey-instrument.pdf} and \texttt{prescreen-} \texttt{survey-instrument.pdf} & 1.7 MB & Questionnaires \\
\bottomrule
\end{tabular}
\end{center}

\noindent Download quirk: \texttt{survey.csv} is ingested by Dataverse as a tab file; the listed md5 refers to the
\emph{original} upload, so the script requests \texttt{?format=original}. Dataverse also returns 403 to requests without a User-Agent.

\subsection*{First look (crude keyword flags, not the audit classification)}
Wearable = title contains Fitbit, Garmin, Apple Watch, smartwatch, fitness/activity tracker or Forerunner. Running =
title/category keywords. The wearable flag includes accessories (watch bands are 43\% of flagged rows), so it overstates device purchases.
Script: \texttt{code/scripts/first\_look\_open\_ecommerce.py}.

\begin{center}\small
\begin{tabular}{lr}\toprule
Item & Value \\\midrule
Purchase rows & 1,850,717 \\
Users with purchases & 5,027 \\
Survey respondents & 5,027 \\
Order dates & 2018-01-01 to 2024-08-15 \\
Median purchases per user & 232 \\
Rows flagged wearable (title keywords) & 6,735 \\
Users with a flagged wearable (ever-treated) & 1,982 \\
  of which $\geq$12 months pre-history & 1,413 \\
  of which $\geq$12 months pre and post & 1,076 \\
Rows flagged running-related & 10,827 \\
Users with a running-related purchase & 2,802 \\
\bottomrule\end{tabular}
\\[1ex]
First flagged wearable purchase, by year:\\[0.5ex]
\begin{tabular}{rrrrrr}\toprule
2018 & 2019 & 2020 & 2021 & 2022 & 2023 \\\midrule
463 & 422 & 359 & 406 & 308 & 24 \\\bottomrule\end{tabular}

\end{center}

\noindent Takeaways for the audit: the data run 2018--Aug 2024; there are roughly 2,000 ever-treated users and about
1,100 with at least a year of history on both sides under the crude flag; cohorts are spread over 2018--2022 (few in 2023).
The category field (e.g.\ \texttt{WEARABLE\_COMPUTER}, \texttt{BIOMETRIC\_MONITOR}, \texttt{SHOES}) should give a cleaner classification than titles.
The survey has no running or exercise items, so the proxy cannot be validated within the dataset.

\section{Packages}
\begin{center}\small
\begin{tabular}{p{4.6cm}p{5cm}p{5cm}}
\toprule
Need & Package & Status \\\midrule
Wrangling large CSV & pandas, polars, pyarrow & already in container \\
Panel/FE, baseline TWFE & linearmodels, statsmodels & already in container \\
Callaway--Sant'Anna (primary) & R \texttt{did} 2.5.1 & added \\
Sun--Abraham, fast FE & R \texttt{fixest} 0.14.2; Python \texttt{pyfixest} 0.60 & added \\
Imputation estimator (robustness) & R \texttt{didimputation} 0.5.1 & added \\
TWFE weight diagnostics & R \texttt{bacondecomp} 0.1.1 & added \\
Sensitivity to parallel-trends violations & R \texttt{HonestDiD} 0.2.8 & added (GitHub, see below) \\
Plots & matplotlib, seaborn, R ggplot2 & already in container \\
Structural extension (if gated in) & pyblp & already in container \\
\bottomrule
\end{tabular}
\end{center}

\noindent Choice: R is the reference implementation for \texttt{did} and \texttt{HonestDiD}, so estimation runs in R inside the
same container, with Python for data prep and a \texttt{pyfixest} cross-check. Everything was added to \texttt{code/Dockerfile};
\texttt{requirements.txt} is empty and was left unchanged.

\noindent Install pitfalls: \texttt{HonestDiD} is not on CRAN (installed from GitHub). Its dependency \texttt{Rglpk} only builds
against conda's \texttt{glpk} (system \texttt{libglpk-dev} is not picked up), and \texttt{CVXR} $\geq$1.1 needs a Rust toolchain,
so \texttt{CVXR} is pinned to 1.0-15. \texttt{arrow} (R) failed to install and is not needed; use \texttt{pyarrow}.

\section{Open items}
\begin{itemize}[nosep]
  \item Docker compose warns that \texttt{version} is obsolete and that an orphan container from another project exists; harmless, not changed.
  \item Mount of \texttt{analysis/} in the compose file is needed to write outputs from the container; it is already present.
  \item Next issue: data audit (clean wearable vs.\ accessory classification, running-gear classification, cohort sizes, pre-trends) and go/no-go.
\end{itemize}
\end{document}
